\section{Further Resources}\label{sec:appendix:further_resources}
\vspace{7.5cm}

Below, we collect further helpful resources for those who wish to deepen their cryptographic knowledge, in particular books (\Cref{sec:appendix:further_resources:books}), online courses (\Cref{sec:appendix:further_resources:courses}), hands-on challenge platforms (\Cref{sec:appendix:further_resources:ctf}), newsletters and podcasts (\Cref{sec:appendix:further_resources:np}), useful websites and reading guides (\Cref{sec:appendix:further_resources:web}), and blogs (\Cref{sec:appendix:further_resources:blogs}).


\subsection{Books}\label{sec:appendix:further_resources:books}

The following books are ordered by latest edition (most recent first):

\begin{itemize}

	\item \textbf{The Joy of Cryptography: An Undergraduate Course in Provable Security} by \textit{Mike Rosulek} (latest edition published on January 2026 with 702 pages): a proof-oriented textbook that introduces the theoretical foundations of modern cryptography. It presents formal security definitions and reduction-based proofs for topics including symmetric and public-key encryption, digital signatures, pseudorandomness, zero-knowledge proofs, and secure multiparty computation, emphasizing rigorous reasoning within the provable security framework. Print book and Ebook can be bought at \url{https://mitpress.mit.edu/9780262049979/the-joy-of-cryptography/} (The MIT Press);
	
	\item \textbf{Serious Cryptography: A Practical Introduction to Modern Encryption} by \textit{Jean-Philippe Aumasson} (latest edition being the 2nd and published on August 2024 with 376 pages): breaking down fundamental math concepts behind ciphers without shying away from how and why they work (and how they can fail), it says to be ``a practical guide to the past, present, and future of cryptographic systems and algorithms''. Print book and Ebook can be bought at \url{https://nostarch.com/seriouscrypto} (No Starch Press);
	
	\item \textbf{Crypto Dictionary: 500 Tasty Tidbits for the Curious Cryptographer} by \textit{Jean-Philippe Aumasson} (latest edition being the 1nd and published on March 2021 with 160 pages): a dictionary (organized alphabetically) with hundreds of cryptographic-related entries and illustrations concerning algorithms, trivia from the history of cryptography, attacks, polemics, the threat of quantum computing, and more. Print book and Ebook can be bought at \url{https://nostarch.com/crypto-dictionary} (No Starch Press);
	
	\item \textbf{Understanding Cryptography: A Textbook for Students and Practitioners} by \textit{Christof Paar, Jan Pelzl, and Tim Guneysu} (latest edition being the 2nd and published on May 2024 with 564 pages): focusing on covering (nearly) all symmetric, asymmetric, and post-quantum cryptographic algorithms, it says to be a ``must-have book [] indispensable as a textbook for graduate and advanced undergraduate courses''. Print book and Ebook can be bought at \url{https://www.cryptography-textbook.com/} (Springer);
	
	\remarkBlock{Video lectures on Understanding Cryptography}{Christof Paar's video lectures --- available at the link \url{https://www.youtube.com/@introductiontocryptography4223/videos} --- are an excellent companion to the Understanding Cryptography book, mixing theory with practical demos.}
			
	\item \textbf{A Graduate Course in Applied Cryptography} by \textit{Dan Boneh and Victor Shoup} (latest edition being the 0.6th and published on January 2023 with 1130 pages): a rigorous, graduate-level book emphasizing provable security, it says to be ``about mathematical modeling and proofs to show that a particular cryptosystem satisfies the security properties attributed to it'' and to ``describe dozens of cryptographic algorithms, give practical advice on how to implement them, and show how they can be used to solve security problems''. PDF freely available at the link \url{https://toc.cryptobook.us/};
	
	\item \textbf{Real-World Cryptography} by \textit{David Wong} (latest edition being the 1st and published on September 2021 with 400 pages): a hands-on book focused on applying cryptographic techniques in real systems, it says to teach ``practical techniques for day-to-day work as a developer, sysadmin, or security practitioner''. Print book, Ebook, and Audiobook can be bought at \url{https://www.manning.com/books/real-world-cryptography} (Manning Publications Co);
	
	\item \textbf{Introduction to Modern Cryptography} by \textit{Jonathan Katz and Yehuda Lindell} (latest edition being the 3rd and published on December 2020 with 648 pages): an academic book offering a rigorous treatment of cryptography from a theoretical standpoint focused on formal security models. PDF can be bought at \url{https://www.taylorfrancis.com/books/mono/10.1201/9781351133036/introduction-modern-cryptography-yehuda-lindell-jonathan-katz} (CRC Press);
	
	\item \textbf{Crypto 101} by \textit{Laurens Van Houtven} (latest edition being incomplete and published on 2017 with 223 pages): an open-source book on cryptography, it says to be ``an introductory course on cryptography, freely available for programmers of all ages and skill levels'' and extensively ``focus on breaking cryptography'', teaching how common flaws can be exploited. PDF freely available at the link \url{https://www.crypto101.io/}.
	
\end{itemize}

\readBlock{Older books}{The following books are perhaps outdated but surely still worth mentioning:
\begin{itemize}
	\item Cryptography Engineering: Design Principles and Practical Applications by Niels Ferguson, Bruce Schneier, and Tadayoshi Kohno (latest edition being the 1st and published on March 2010 with 384 pages: Print book and Ebook can be bought at \url{https://www.schneier.com/books/cryptography-engineering/} (Wiley);
	\item Handbook of Applied Cryptography by Alfred J. Menezes, Paul C. van Oorschot, and Scott A. Vanstone (latest edition being the 5th and published on August 2001 with 810 pages --- latest errata update on 2014): PDF freely available at the link \url{https://cacr.uwaterloo.ca/hac/} (CRC Press);
	\item Applied Cryptography by Bruce Schneier (latest edition being the 2nd and published on January 1996 with 784 pages): Print book and Ebook can be bought at \url{https://www.schneier.com/books/applied-cryptography/} (John Wiley \& Sons).
\end{itemize}}


\subsection{Online Courses}\label{sec:appendix:further_resources:courses}

The following online courses are ordered alphabetically (note that we include only courses available to everyone; university courses targeting students are not included):

\begin{itemize}
	
	\item \textbf{Cryptography} by \textit{University of Maryland, Jonathan Katz}: an online course that explains the basis of traditional cryptography. Freely available at the link \url{https://www.coursera.org/learn/cryptography};
	
	\item \textbf{Cryptography I} by \textit{Stanford University, Dan Boneh}: a very popular online course that provides a broad introduction to cryptographic primitives and how to use them correctly. The course starts from basic symmetric ciphers and works up to public-key cryptography. Later modules examine real protocols and common implementation mistakes, and an optional programming project lets students experiment with breaking a toy system. Freely available at the link \url{https://www.coursera.org/learn/crypto/}.
	
	\warningBlock{Cryptography II was removed}{The course Cryptography I used to have a follow up named \textbf{Cryptography II} covering more advanced topics like \glspl{zkp}. However, Cryptography II was recently removed from the online course platform.\footnote{Cryptography II by Dan Boneh has been removed from Coursera, any alternative sources? : r/crypto (\url{https://www.reddit.com/r/crypto/comments/1deuruf/cryptography_ii_by_dan_boneh_has_been_removed/})}}
	
\end{itemize}


\subsection{Hands-On Challenge Platforms}\label{sec:appendix:further_resources:ctf}

The following hands-on challenge platforms are ordered alphabetically:

\begin{itemize}
	\item \textbf{CryptoHack}: a gamified platform for learning cryptography by breaking bad implementations of \gls{aes}, \gls{rsa}, and \gls{ecc}. The platform is organized as \gls{ctf}-style challenges, escalating from \texttt{xor} ciphers to lattices. Notably, the platform is active and regularly updated with new challenges. Freely available at the link \url{https://cryptohack.org/};
	
	\item \textbf{Cryptopals Crypto Challenges}: a famous (fixed) set of programming challenges to demonstrate real-world attacks on cryptographic algorithms (not side-channel attacks, though), spanning from \texttt{xor} ciphers and \gls{aes} in \gls{ecb} mode up to \gls{rsa} and \gls{ecc}. Freely available at the link \url{https://cryptopals.com/};
	
\end{itemize}

\readBlock{Other Platforms}{The following platforms are perhaps smaller in scope but surely still worth mentioning:
	\begin{itemize}
		\item id0-rsa, freely available at the link \url{https://id0-rsa.pub/};
		\item OverTheWire community's Krypton, freely available at the link \url{https://overthewire.org/wargames/krypton/};
		\item HackTheBox Crypto challenges, , freely available at the link \url{https://help.hackthebox.com/en/articles/5185436-how-to-play-challenges}.
\end{itemize}}


\subsection{Newsletters and Podcasts}\label{sec:appendix:further_resources:np} 

The following newsletters and podcasts are ordered alphabetically (newsletters first):

\begin{itemize}

	\item \textbf{Crypto-Gram}: a ``free monthly newsletter providing summaries, analyses, insights, and commentaries on security technology'' by Bruce Schneier. Despite the name, Crypto-Gram is not only about cryptography, but instead it covers broader cybersecurity. Crypto-Gram often highlights recent cryptographic breaches or debates (e.g., encryption backdoors, algorithm policy) with a balanced and clear analysis. Freely available at the link \url{https://www.schneier.com/crypto-gram/};
	
	\item \textbf{Cryptography Dispatches}: a newsletter by Filippo Valsorda where each issue is a short essay focusing on a specific topic of applied cryptography or recent development. Started around 2020, this newsletter is valuable for providing digestible ``dispatches'' from the world of applied cryptography in plain language. Freely available at the link \url{https://buttondown.com/cryptography-dispatches/};

	\item \textbf{Cryptography FM}: a podcast hosted by Nadim Kobeissi and Lúcás Meier offering ``in-depth, substantive discussions on the latest news and research in applied cryptography''. The latest episode dates back to February 2023. Freely available at the link \url{https://cryptography.fireside.fm/};

	\item \textbf{Security Cryptography Whatever}: a podcast hosted by David Adrian, Deirdre Connolly, and Thomas Ptacek where ``some cryptography and security people talk about security, cryptography, and whatever else is happening''. Episodes range from deep dives into specific protocols or vulnerabilities to casual chats about industry news. The podcast is active as of 2025. Freely available at the link \url{https://securitycryptographywhatever.com/}.
	
\end{itemize}


\subsection{Useful Websites and Reading Guides}\label{sec:appendix:further_resources:web}

The following websites and guides are ordered alphabetically:

\begin{itemize}
	\item \textbf{Crypto Gotchas!}: a curated list of common mistakes and pitfalls in cryptography by Greg Rubin. The list enumerates concerns like ``keys wearing out'', poor entropy, and padding oracles attacks. Essentially, Crypto Gotchas! is a checklist of what not to do in applied cryptography, with links to blog posts or papers. Freely available at the link \url{https://gotchas.salusa.dev/};
	
	\item \textbf{r/crypto}: a cryptography community/wiki on Reddit that serves as a beginner's guide, FAQ, and Q\&A. Freely available at the link \url{https://www.reddit.com/r/crypto/wiki/index/};
	
	\item \textbf{The Stick-Figure Guide to \gls{aes}}: a famous illustrated blog post published in 2009 by Jeff Moser that provides a ``layman's introduction to cryptography and AES''. The blog post uses intuitive analogies and even stick-figure cartoons to explain how \gls{aes} works internally, allowing for understanding \gls{aes} encryption at a conceptual level (e.g., what are S-boxes and mix-columns) without too many mathematical details. Freely available at the link \url{http://www.moserware.com/2009/09/stick-figure-guide-to-advanced.html};
	
	\item \textbf{Timeline of Cryptography}: a Wikipedia article that provides a chronological timeline of important events in the history of cryptography, from classical ciphers in Roman era up to recent developments. Freely available at the link \url{https://en.wikipedia.org/wiki/Timeline_of_cryptography}.
\end{itemize}

	
\subsection{Blogs}\label{sec:appendix:further_resources:blogs}

The following blogs are ordered alphabetically:

\begin{itemize}
	
	\item \textbf{A Few Thoughts on Cryptographic Engineering}: a blog by prof. Matthew Green containing ``some random thoughts about crypto. Notes from a course I teach. Pictures of my dachshunds''. This blog is among the most respected for applied cryptography, offering long-form explanations of things like how Apple's iPhone encryption works, analyses of protocol designs, critiques of government policies, and timely explainers of new vulnerabilities. Freely available at the link \url{https://blog.cryptographyengineering.com/};
	
	\item \textbf{Chosen Plaintext's blog}: a blog by Chosen Plaintext Consulting (a security consulting company) which features deeply technical articles on cryptography engineering. Many of the posts focus on secure implementation and cryptographic pitfalls (e.g., ``A beginner's guide to constant-time cryptography'' and ``Cryptography Red Flags''). Almost all posts date back to 2015. Freely available at the link \url{https://www.chosenplaintext.ca/blog/};
	
	\item \textbf{Dhole Moments}: a blog by Soatok covering software, security, and cryptography from the perspective of an independent researcher. The content is practitioner-oriented; for example, Soatok has articles debunking bad cryptographic advice, explaining secure protocol design, and reviewing the cryptography in messaging apps. Freely available at the link \url{https://soatok.blog/author/soatok/};
	
	\item \textbf{Kelby Ludwig's blog}: a blog containing technical posts on specific cryptographic topics suc as \gls{ecc}, the Hidden Number Problem, \gls{ecdsa} quirks, and lattice-based cryptography. The latest post dates back to March 2021. The blog is not a general beginner blog, but it is great for niche topics. Freely available at the link \url{https://kel.bz/};
	
	\item \textbf{Latacora's Blog}: a blog by Latacora (a security consulting company) giving practical guidance on cryptography for developers, especially for what concerns the selection of cryptographic algorithms, secure defaults, libraries, and post-quantum cryptography as well. Freely available at the link \url{https://www.latacora.com/blog/};
	
	\item \textbf{Neil Madden's blog}: ``Thoughts on application security, applied crypto, philosophy and logic'', often deep diving into cryptographic implementations, protocols, and vulnerabilities in an approachable style --- especially for developers. Freely available at the link \url{https://neilmadden.blog/};
	
	\item \textbf{Ryan Castellucci's blog}: ``posts on computer security, programming, systems administration, electronics, and general geekery''. The blog has not been frequently updated lately, but the existing content is insightful, especially on topics of practical cryptanalysis. Freely available at the link \url{https://rya.nc/}.
	
\end{itemize}

\readBlock{Other Blogs}{David Wong also used to curate a large list of blogs related to cryptography --- available at the link \url{https://github.com/mimoo/crypto_blogs} --- although the latest update to the list dates back to 2021.}
